\section{The remaining global gap and the scope of the methods}
\label{sec:global}

The finite and exterior theorems give the coverage uniform over the
entire time slab in Table~\ref{tab:coverage}. The endpoints $1001$ and
$e^{2000}$ are included in their respective proved regions.
Section~\ref{sec:uniform-edge} also records a small additional box near
$4\pi e^{20}$ on the thin time interval $[1/20-10^{-14},1/20]$.
That island does not close the bridge over the stated slab.
\begin{table}[htbp]
\centering
\caption{Real first-jet coverage for $1/100\le t\le1/5$.}
\label{tab:coverage}
\begin{tabular}{@{}ll@{}}
\toprule
Spectral region & Status\\
\midrule
$|x|\le1001$ & Theorem~\ref{thm:finite-height}\\
$1001<|x|<e^{2000}$ & Bridge not covered over the full slab\\
$|x|\ge e^{2000}$ & Theorem~\ref{thm:outer20}\\
\bottomrule
\end{tabular}
\end{table}

The published comparison baseline is $\Lambda\le1/5$
\cite[Corollary 2]{PlattTrudgian}. No improved Newman bound follows
from the two disconnected covered regions. In particular, approximation
accuracy relative to $\mathcal G_0$ is not a lower bound on the
complete field or its first jet.

\subsection{Limits of the present explicit-mode implementation}

For either contour rule in Section~\ref{sec:effective}, the retained
list has
\[
 N(x)=\lceil\sqrt{bx}\rceil-1,\qquad
 \sqrt{bx}-1\le N(x)<\sqrt{bx}.
\]
The phase-gap premise itself forces $m^2>x/(4\pi)$.
At $b=1/8$ and $x=e^{2000}$, direct outward evaluation gives
$\log_{10}N(x)>433.8429369097558$.
The present implementation materializes every retained mode and
visits each at every quadrature callback. One such callback thus
already entails more than $10^{433}$ mode operations.
This excludes that enumeration scheme as a feasible bridge to
the present onset, not faster theta assembly or other representations.

The retained source interval must also grow under the current
absolute tangent payment. Its weakest $n=1$, $j=Y=0$ condition
requires $4c_d\pi e^{4U}>9$. For the constant-distance rule,
$c_d\le4A_c/x$, and therefore
\begin{equation}\label{eq:source-growth}
 U>\frac{\log x}{4}+
       \frac14\log\frac9{16\pi A_c}.
\end{equation}
At $x=e^{2000}$, $A_c=33$, this requires $U>498.69584960$.
For logarithmic distance with $a=4$ the corresponding necessary
bound is $U>\log x/4+\log(9/(64\pi\log x))/4$, exceeding
$497.32317728$ there. These are conditions on the current error
payment, not lower bounds on the true oscillatory tail.
At $U=2$, $T=1/5$, $Y=1$, the logarithmic-rule denominator obeys
\[
 D_{1,0}\le64\pi e^8\frac{\log x}{x}-10.8<0
 \quad(x\ge10^6).
\]
For constant $A_c\le33$ it is at most
$16\pi A_c e^8/x-10.8<0$ on the same range. Thus the fixed
$U=2$ contract cannot simply be reused at larger height.

There is a separate restriction on the fixed-degree majorant.
Assume the constant-distance rule, $U\ge1/4$,
$A_cb\ge1/(\pi e)$, $bx\ge4$, and $x\ge64\pi A_ce$.
Let $R_0=\sqrt{x/(4\pi A_ce)}\ge4$.
All integers $n\le\lfloor R_0\rfloor$ are retained because
$R_0\le\sqrt{bx}/2\le m-1$.
For $0\le u\le1/4$,
$\pi c_dn^2e^{4u}\le1$ for these modes.
Their upper half gives $\sum n^4\ge R_0^5/1024$.
Consequently the normalized positive common moment of any order $J$
satisfies
\begin{equation}\label{eq:moment-growth}
 M_J\ge
 \frac{\pi^2e^{A_c-1}}{2048(4\pi A_ce)^{5/2}}x^{3/4}.
\end{equation}
This follows by retaining their $2\pi^2n^4$ source coefficients
over that interval; the time/source factors and
$1+(u+d)^J$ are at least one.
At the wall endpoint $y=1$, the degree-24 Taylor payment about
the real midpoint contains $M_{52}/25!$. It eventually exceeds
any fixed $\mathcal G_0$-relative allowance such as $10^{-8}$.
Subdividing cells without changing that midpoint or those degrees
does not change this bound. Raising degrees, recentering, using a
different majorant, or allowing a growing independently proved
physical margin are separate possibilities not excluded here.

\subsection{The alternative complex barrier contract}

For precision, \cite[Theorem 1.2]{Polymath} uses $t_0,X>0$ and
$0<y_0\le1$ with three hypotheses:
\begin{enumerate}
\item no initial zeta zeros with $(1+y_0)/2\le\Re s\le1$
and $0\le\Im s\le X/2$;
\item no zeros of $H(t_0,x+iy)$ with
$x\ge X+\sqrt{1-y_0^2}$ and
$y_0\le y\le\sqrt{1-2t_0}$;
\item no zeros for $0\le t\le t_0$ with
$X\le x\le X+\sqrt{1-y_0^2}$ and
$\sqrt{y_0^2+2(t_0-t)}\le y\le\sqrt{1-2t}$.
\end{enumerate}
Its conclusion is $\Lambda\le t_0+y_0^2/2$.
Only parameter choices with the indicated real boundary heights
are used here.

Choose $0<t_0\le1/5$ and $X=1000-\sqrt{1-y_0^2}$.
Initial verification follows from \cite[Theorem 1]{PlattTrudgian},
using a height below $500$. The complete moving barrier lies in
$999<x\le1000$ and above the real axis, so
Lemma~\ref{lem:finite-propagation} and its theta application supply
the third hypothesis. Proposition~\ref{prop:wall} pays the
final seam at $x=1000$. The outstanding hypothesis is the
entire final-time complex exterior for $x>1000$.

For example, $t_0=1/10$, $y_0=2/5$ would yield
$\Lambda\le9/50=0.18$ if one additionally proved nonvanishing on
\[
 t=1/10,\qquad x>1000,\qquad 2/5\le y\le\sqrt{4/5}.
\]
This is a conditional contract, not an obtained bound.
The real first-jet theorem \ref{thm:outer20} is not this complex
nonvanishing assertion. Nor can a new finite wall silently extend
the range of the initial zeta verification it would require.
Reusing the known all-real time $t_0=1/5$ gives only
$1/5+y_0^2/2>1/5$; making the exterior empty by taking
$y_0>\sqrt{1-2t_0}$ gives a bound above $1/2$.

\subsection{A remote polynomial control for the finite premises}

The next example isolates a limitation of the finite-wall, strip and
all-real-baseline premises. It does not have the theta square source,
and is not a counterexample to any theta theorem.

\begin{proposition}[A collision inside the uncovered bridge]
\label{prop:remote-control}
Let $L=2000$, $q=2/5$, and define
\[
 F_t(z)=z^4-2A(t)z^2+B(t),\qquad
 A(t)=L^2-q+6t,\quad
 B(t)=(L^2+q)^2+4t(L^2-q)+12t^2.
\]
Then $F_t$ satisfies the backward heat equation, evenness and
real conjugation symmetry. It has no zeros on the complete wall
$\Re z=1000$ or in $|\Re z|\le1001$ for $0\le t\le1/5$,
is positive on the imaginary axis, satisfies the contracted unit
zero strip, and has all zeros real at $t=1/5$.
Nevertheless its all-real threshold $\lambda_L$ satisfies
\[
 0.1999999875<\lambda_L<0.19999999375,
\]
with real double zeros at that time. It also has no real
first-jet collision beyond the exterior in Theorem~\ref{thm:outer20}.
\end{proposition}
\begin{proof}
Coefficient differentiation gives $F_t=-F_{zz}$, and
$F_0(z)=((z-L)^2+q)((z+L)^2+q)$.
Put $f(t)=6t^2+2(L^2-q)t-L^2q$. Then $A^2-B=4f(t)$.
The function $f$ increases for $t\ge0$, with positive root
\[
 \lambda_L=
 \frac{\sqrt{(L^2-q)^2+6L^2q}-(L^2-q)}6.
\]
Before this time the roots are two nonreal conjugate pairs;
at it they are two real double roots; afterwards all four are real.
Exact rational evaluations of $f$ at the displayed decimal endpoints
give the strict bracket. More generally
$1/5-1/L^2<\lambda_L<1/5$.

For nonreal roots, their real parts squared equal
$(\sqrt B+A)/2\ge A(0)/2>1400^2$.
For real roots before $t=1/5$,
$4f(t)\le4f(1/5)=8/25<1$, so their squared moduli are at
least $A(0)-1>1400^2$. All roots therefore remain beyond
$|\Re z|>1400$. In particular
$|F_t(1000+iy)|>400^4$ and
$|F_t(x)|>399^4$ for $|x|\le1001$.
Also $F_t(iy)=y^4+2A(t)y^2+B(t)>0$.

For a nonreal root its imaginary part squared is
$(\sqrt B-A)/2$. Since
$(L^2+q+2t)^2-B(t)=8t(q-t)\ge0$ up to the threshold,
this is at most $q-2t\le1-2t$. After the threshold
all roots are real. Initial nonreal roots lie outside the
finite verified region, so the initial-reality hypothesis there
is also satisfied.

Finally every root stays in $|z|<3000$ through $0\le t\le1/2$.
For nonreal roots use $|z|^4=B(t)\le B(1/2)<3000^4$;
for real roots use
$z^2\le A(1/2)+\sqrt{4f(1/2)}<A(1/2)+4000<3000^2$.
The collisions therefore lie below
$e^{20/\varepsilon}\ge e^{40}>3000$.
\end{proof}

For varying $L$, the finite-wall, strip and baseline premises admit
thresholds tending to $1/5$. The new exterior is not imposed uniformly
on that entire varying family. For the fixed $L=2000$ example,
even adding the exterior conclusion leaves a collision in the bridge.
The example fails the stronger inequality $\Lap_0>E/100$ at
$(t,x)=(1/100,1000)$, as exact rational substitution verifies.
Thus it supplies no obstruction to an argument that uses that
additional signed information outside its presently certified rectangle.

The next missing input is a quantitative estimate for the retained
complete arithmetic sum, a feasible complete real bridge, or the
specified final-time complex exterior. Full modular reflection
identifies the theta source in the coefficient class of
Theorem~\ref{thm:rigidity}; converting that rigidity into such an
estimate remains open. All-height noncollision is not established
by giving another name to that missing estimate.
